\documentclass[a4paper]{article}
\usepackage{luatexja}
\usepackage{amsmath}
\usepackage{amssymb}
\usepackage{geometry}
\usepackage{hyperref}

\geometry{margin=25mm}

\title{暗黒物質を要しない銀河回転曲線：\\
双対時空のコンパクト性}

\author{https://github.com/hypernumbernet}
\date{\today}

\begin{document}

\maketitle

\begin{abstract}
核と電子殻を層状にする同一の双対ローター幾何は、宇宙論的スケールではコンパクト性として展開する。双対ローターは三次元球面である。慣性質量に結びついた $(\phi,\theta)$ が駆動する重力チャンネルを、コンパクト化半径 $R$ の銀河スケールのパッチとして、この球面を組織するものと読む。電荷に結びついた $(\alpha,\beta)$ が駆動する電磁チャンネルは、実質的により大きな、ほぼ宇宙的なスケールにわたって展開する。三次元球面上の球対称な源は、ガウスの法則で正規化された、半径方向に調和な余接ポテンシャルをもつ。その加速度は、周期 $\pi R$ で繰り返される源の像についての和である。像を通常の半径に沿って並べると、対蹠点 $r=\pi R$ で円速度は有限な値をもたず、天の川型の質量ではその対蹠点が測定された回転曲線の範囲に入る。その配置は通常の半径を周期的にしてしまうが、コンパクトなのは双対時空だけである。同じ像を双対方向に沿って並べれば、加速度は $(GM/(rR))\coth(r/R)$ となる。コンパクト化半径が $R=\sqrt{GM/a_0}$ とスケールするなら、これは観測加速度とニュートン加速度の間のただ一つの関係であり、回転曲線は、バリオンのニュートン曲線に $x\coth x$（$x^2=a_0r/v_{\text{bar}}^2$）を掛けたものになる。曲線は中心から立ち上がり、折れ曲がり、再び上昇することなく平坦になる。その速度は $v^4=GMa_0$ を厳密に満たす。銀河がそのようなパッチであるなら、不足する加速度はこの幾何であり、バリオンだけから読まれる。その同定が、提案する読みである。以下の加速度 $a_0$ は外部の入力であり、双対ローター代数の帰結ではない。
\end{abstract}

\section{はじめに}
\label{sec:introduction}

二重時空理論は共有連続体を拒む。各質量粒子は通常のミンコフスキー枠と双対枠を運び、重力はその対のねじれ不整合である~\cite{dst-main}。双対ローターは三次元球面である。銀河の尺度では、慣性質量に結びついた $(\phi,\theta)$ が駆動する重力チャンネルを、コンパクト化半径 $R$ のパッチとしてこの球面を組織するものと読む。電荷に結びついた $(\alpha,\beta)$ が駆動する電磁チャンネルは、より大きな半径の上で展開する。$R=\sqrt{GM/a_0}$ を定める加速度 $a_0$ と、のちに電磁半径として読むハッブル間隔は、外部の入力である。どちらも双対ローター代数の帰結ではない。

第\ref{sec:compact}節は三次元球面とその等距離図法を記録する。第\ref{sec:cotangent}節は余接ポテンシャルを書く。第\ref{sec:radial}節は像を通常の半径に沿って並べ、その配置が銀河の法則ではないことを示す。第\ref{sec:dual}節は同じ像を双対方向に沿って並べる。第\ref{subsec:btfr}節はパッチを質量とともにスケールし、加速度関係を得る。第\ref{sec:rotation}節はバリオンのニュートン曲線から回転曲線を作る。第\ref{sec:sky}節は、天空が決めうるものを切り出す。

\section{コンパクトな双対時空}
\label{sec:compact}

双対ローター $R_{\text{dual}} \in \mathrm{SU}(2) \cong S^3$ は、すべての質量粒子に内在する双対自由度をパラメトライズする。三つの回転角は半径 $R$ の三次元球面上の座標であり、通常時空は平坦なままである。これらの角は座標としての読みである。各粒子は三次元球面の元を担い、粒子 $b$ の粒子 $a$ に対する位置は、一方の元を他方へ送るローター $\rho(a,b)=a^{-1}b$ である。このローターが恒等元であるのは二つの元が一致するときに限り、それらのなす角の余弦は $\rho$ の実部である。三次元球面の回転は、あらゆる元に一度に作用する単位四元数の対である。左の因子は $\rho$ を変えない。右の因子は $\rho$ を $h^{-1}\rho h$ に置き換え、角は変わらない。等距離図法の極を $g$ へ移すと、元 $p$ の座標は $g^{-1}p$ の対数になる。座標は変わり、$\rho$ は変わらない。極を $a$ に置くと、$b$ の座標は $\rho(a,b)$ の対数である。このローターは二つの粒子の関係であり、一つの粒子の通常ローターと双対ローターの不整合ではない。対数写像
\[
\log : S^3 \to \mathfrak{su}(2) \cong \mathbb{R}^3
\]
は各ローターを、半径 $\pi R$ のユークリッド球の一点へ送る。これは三次元の方位正距図法である。原点では忠実であり、対蹠点では特異であって、その一点が球の境界に沿って広がる。体積の歪みはヤコビアン
\[
\mathcal{J}(\chi) = \frac{\sin^2 \chi}{\chi^2}, \qquad \chi = r/R
\]
である。$\chi \to 0^+$ では、$\mathcal{J}$ は $1$ に近づく。対蹠点で $\mathcal{J}$ は消える。ニュートン計算は、あらゆる半径で $\mathcal{J}=1$ と仮定する。銀河スケールのパッチでは、その仮定が成り立つのは中心の近くに限る。距離 $\pi$ のあらゆる方向は同じ一つの元へ持ち上がる。極に対する相対ローターは $-1$ であり、角は $\pi$ である。角が $\pi$ になるのはこの対蹠点に限る。図法の式は、この元も極も原点に書く。$-1$ は回転群の中心にあり、球面のあらゆる回転がこの相対ローターを固定する。一般の相対ローターは右の因子によって共役を受ける。$-1$ はその共役で動かない。図が極と対蹠点を一つの原点に重ねるとき、位置を保つのはこのローターである。この元 $-1$ は、$J_\varepsilon$ の指数がパラメーター $\pi/\varepsilon$ でとる値である。原点を除けばヤコビアン $\mathcal{J}$ は $\sin\chi/\chi$ の平方であり、したがって対蹠点で消える。$\chi=\varepsilon t$ に沿って $\varepsilon\to 0$ とすると、この因子はあらゆる有限の $t$ で $1$ に向かい、指数は親理論の零回転へ向かう~\cite{dst-main}。対蹠点は遠ざかり、平坦な図が保たない方向は零直線である。

\section{余接ポテンシャル}
\label{sec:cotangent}

三次元球面の測地極座標
\[
ds^2 = dr^2 + R^2 \sin^2(r/R)\,(d\vartheta^2 + \sin^2\!\vartheta\, d\varphi^2)
\]
において、測地球面の面積は $4\pi R^2\sin^2(r/R)$ である。余接
\[
\Phi(r) = -\frac{GM}{R}\cot\!\left(\frac{r}{R}\right)
\]
は、そのようなあらゆる球面を通る流束が $-4\pi GM$ となる動径ポテンシャルである。原点の外では、面積因子と $\partial_r\Phi$ の積は定数 $GM$ であり、したがって半径方向のラプラシアンは $\Phi$ を消す。
\[
\partial_r^2\Phi + \frac{2}{R}\cot\!\left(\frac{r}{R}\right)\,\partial_r\Phi = 0.
\]
微分すれば加速度は
\[
g(r) = -\frac{d\Phi}{dr} = -\frac{GM}{R^2\sin^2(r/R)}
\]
である。開区間 $0<r<\pi R$ の上で、この加速度は厳密に負である。力は赤道のどちら側でも質量を向く。ポテンシャルそのものが負であるのは $r<\pi R/2$ のときに限り、消えるのは赤道においてのみである。赤道は負の井戸の端である。力の反転ではない。それを越えると $\Phi$ は正であり、それでも源に向かって下る。測地球面が対蹠点へ向かって縮むにつれ、同じ流束はより狭い面積へ押し込められ、$|g|$ は増大するが、向きは内向きのままである。

$0<r\le R$ では、原点の芽はニュートン的である。ポテンシャルは $-GM/r + GMr/(3R^2)$ と、次数 $r^3/R^4$ の剰余だけ異なる。補正 $GMr/(3R^2)$ は三次元球面の曲率である。非コンパクト化 $R\to\infty$ でそれは消え、ニュートンポテンシャルからの逸脱もともに消える。

\section{半径に沿う像}
\label{sec:radial}

加速度は格子和である。余接の部分分数展開により
\[
\frac{GM}{R^2\sin^2(r/R)}=\sum_{n\in\mathbb{Z}}\frac{GM}{(r-n\pi R)^2}
\]
となる。$n=0$ の項はニュートンの項であり、残りは、半径に沿って $\pi R$ の倍数ごとに再び見える源である。平坦チャート上の円軌道は
\begin{equation}
\label{eq:s3-circular}
v^2(r) = r\,|g(r)| = \frac{GM}{R}\,f(r/R), \qquad f(x)=\frac{x}{\sin^2 x}
\end{equation}
を満たす。区間 $(0,\pi/2)$ の上で形状はただ一つの浅い極小をもち、それは $\tan x=2x$ の根 $x_0\in(1.14,1.20)$ にあって、$1.35<f(x_0)<1.41$ である。その先で $f$ は上昇し、$x\to\pi^-$ で $+\infty$ に向かう。像 $n=1$ が対蹠点で軌道に達するからである。天の川型の質量では、第\ref{subsec:btfr}節のスケーリングのもとで対蹠点は $25$ から $29\,\mathrm{kpc}$ の間にあり、測定された回転曲線はそのような半径を越えて平坦なままである。半径方向の格子は通常の半径そのものを周期的にするが、コンパクトなのは双対時空だけである。これは中心の近くで忠実なチャートによる像の読みであり、銀河の法則ではない。

\section{双対方向に沿う像}
\label{sec:dual}

双対性 $X\mapsto Xi$ は、双対セクターの楕円幾何と通常セクターの双曲幾何とを親理論の解析接続によって結ぶ~\cite{dst-main}。余接の上では、それは $i\cot(ix)=\coth x$ として働く。その接続を通すと、周期 $\pi R$ は通常の半径を離れ、双対方向に沿って走る。像 $n$ は二乗距離 $r^2+(n\pi R)^2$ にあり、一辺は通常空間に、もう一辺はコンパクトな双対円の周りにある。通常の半径は非コンパクトのままである。同じ余接展開を虚の引数で評価すると、この格子は閉じた形に和がとれる。
\begin{equation}
\label{eq:dual-lattice}
g(r)=\sum_{n\in\mathbb{Z}}\frac{GM}{r^2+(n\pi R)^2}=\frac{GM}{rR}\coth\frac{r}{R}.
\end{equation}
ニュートン値に対する比は双対増強
\[
\eta(x)=x\coth x, \qquad x=r/R
\]
であり、あらゆる $x>0$ に対して
\[
\max(1,x)<\eta(x)<\min\!\left(1+x,\;1+\frac{x^2}{3}\right)
\]
が成り立つ。増強は狭義に増大し、原点で $1$ に近づく。一方 $\coth$ は狭義に減少して $1$ に近づく。中心の近くで二つの格子は一致する。どちらも、余接の芽の曲率補正を含む $GM/r^2+GM/(3R^2)$ から始まる。双対和はあらゆる半径でその芽より厳密に小さく、$GM/r^2$ と $GM/(rR)$ のどちらよりも厳密に大きい。大きな半径では後者が支配する。加速度は $1/r^2$ ではなく $1/r$ で落ち、特異な半径はない。

コンパクト化半径を固定すれば、格子は源の質量について線形である。質量の和の加速度は加速度の和であり、質量を定数倍すれば加速度も同じ因子で倍になる。逆二乗則に対する割合の超過は $r/R$ だけに依る。それは正であり、$(r/R)^2/3$ より厳密に小さい。区間 $(0,\pi)$ では、体積ヤコビアンの逆数である三次元球面の因子 $(x/\sin x)^2$ が共有の芽 $1+x^2/3$ を厳密に上回り、双対増強はその芽を厳密に下回る。同じ外部入力のもとで、太陽質量のコンパクト化半径は $7020$ から $7040$ 天文単位の間にあり、したがって一天文単位での双対の割合の超過は $10^{-8}$ を下回る。天の川のコンパクト化半径では、外側ラピディティがニュートン値を上回る割合は $10^{-6}$ を下回り、その傍らで $r=R$ での双対の超過は $0.31$ から $0.32$ の間にある。

\section{バリオン的タリー--フィッシャー関係と加速度関係}
\label{subsec:btfr}

コンパクト化半径が質量とともに $R=\sqrt{GM/a_0}$ とスケールするとする。ここで $a_0 \approx 1.2 \times 10^{-10}\;\mathrm{m/s^2}$ は観測されるバリオン的タリー--フィッシャー加速度であり、双対ローター代数から導かれた量ではなく入力として置く。より重い銀河はより大きなパッチである。$a_0$ の導出は主張しない。点質量の周りの円速度は $v^2=r\,g=(GM/R)\coth(r/R)$ である。それは半径とともに狭義に減少して
\[
v_\infty^2=\frac{GM}{R}, \qquad v_\infty^4=GM\,a_0
\]
に近づく。これは厳密であり、傾きは四で、形状因子はない。接近は上から、しかも速い。$r=3R$ を越えると $v^2$ は $v_\infty^2$ の $1/200$ 以内にあり、したがって $v$ は $v_\infty$ の $0.25$ パーセント以内にある。二度目の上昇も、対蹠点もない。

同じスケーリングは、格子を動径加速度関係に変える。$g_{\text{N}}=GM/r^2$ とおけば $x^2g_{\text{N}}=a_0$ かつ $GM/(rR)=\sqrt{a_0g_{\text{N}}}$ であり、加速度は $g_{\text{N}}$ だけの関数になる。
\begin{equation}
\label{eq:dual-rar}
g=\sqrt{a_0\,g_{\text{N}}}\,\coth\sqrt{\frac{a_0}{g_{\text{N}}}}.
\end{equation}
その範囲は
\[
\max\!\left(g_{\text{N}},\sqrt{a_0g_{\text{N}}}\right)<g<g_{\text{N}}+\min\!\left(\sqrt{a_0g_{\text{N}}},\,\frac{a_0}{3}\right)
\]
である。高加速度ではニュートン値からの超過は $a_0/3$ を下回る。パッチ半径を囲まれた質量に結ぶと、$k>1$ に対して $g_{\text{N}}$ を $k$ 倍した加速度は、もとの加速度の $k$ 倍より厳密に小さく、二つのニュートン寄与を合わせた加速度はそれぞれの加速度の和より厳密に小さい。$g_{\text{N}}$ が $a_0$ に比べて大きいあいだは逆二乗の読みが加速度を運び、$g_{\text{N}}$ が $a_0$ を通過するにつれて割合の超過は一のオーダーになる。低加速度では $g$ は $\sqrt{a_0g_{\text{N}}}$ に近づく。ミルグロムの現象論の深い領域である。その間で関係は確定しており、選ぶべき補間関数はない。$g_{\text{N}}=a_0$、すなわち $x=1$ では、増強は $\coth 1$ であり、$1.31$ と $1.32$ の間にある。

天の川の代入質量 $M \approx 6 \times 10^{10}\,M_\odot$ に対し、同じ入力は $8\,\mathrm{kpc}<R<9\,\mathrm{kpc}$ と
\[
175\,\mathrm{km/s}<v_\infty<176\,\mathrm{km/s}
\]
を与える。$a_0$ をパッチの特性加速度と同定することは、なお提案である。

\section{回転曲線}
\label{sec:rotation}

半径 $r$ における流束は囲まれた質量 $M(r)$ のものであり、それを担うパッチは $R(r)^2=GM(r)/a_0$ をもつ。バリオンの外では、これは全バリオン質量 $M_{\text{b}}$ のただ一つのパッチである。内側では、囲まれた質量とともに縮む。$g_{\text{N}}=v_{\text{bar}}^2/r$ とおけば、軌道は $x=r/R(r)=\sqrt{a_0/g_{\text{N}}}$ にあり、回転曲線は
\begin{equation}
\label{eq:dual-rotation}
v^2(r)=v_{\text{bar}}^2(r)\,x\coth x, \qquad x^2=\frac{a_0\,r}{v_{\text{bar}}^2(r)}
\end{equation}
となる。これは $g_{\text{N}}=v_{\text{bar}}^2/r$ における関係 \eqref{eq:dual-rar} の $r$ 倍である。ここで $v_{\text{bar}}^2=v_\star^2+v_{\text{gas}}^2+v_{\text{bulge}}^2$ は、ニュートン的な質量模型とまったく同じく、測光と 21-cm 地図から計算する。調整できる量は恒星の質量光度比だけである。$a_0$ は普遍であり、ハローはない。球対称な質量では、この処方は厳密である。円盤では、円盤のニュートン曲線に、半径ごとに同じ関係を当てはめる。二つの場合は、銀河が育つ過程でつながる。始まりは、ランダムに運動する恒星の集合体である。それらの角運動量は打ち消し合い、質量は球である。引力は三次元球面の大円に沿う。移動は、その球面の方位等距離図法である平坦なチャートの、直線の弦である。チャートを平坦として扱うのは、銀河が宇宙の大きさの三次元球面に比べて小さいからにすぎず、厳密には移動もその巨大な三次元球面上にある。体をチャート距離 $\chi$ に、源をチャート距離 $\psi$ に置き、二つの動径のなす角を $\alpha$ とする。大円の初期接線をチャートで読むと、源のまわりのモーメントは
\begin{equation}
\label{eq:sphere-chord}
\left(\frac{\chi^{2}\sin\psi}{\sin\chi}-\frac{\chi\psi\sin\psi}{\sin\chi}\cos\alpha-\psi\sin\chi\cos\psi+\psi\sin\psi\cos\alpha\cos\chi\right)\sin\alpha
\end{equation}
であり、二つの動径が張る平面に垂直である。体と源がチャート原点からの同一半直線上にあれば、モーメントは消える。源が原点そのものであれば、やはり消える。原点を通る大円は直線の半直線であり、引力は平坦なチャートで中心方向のままである。このモーメントはその地図における読みである。体と源と極を同じ左回転で運べば、モーメントは変わらない。一つの対では、極が恒等元のときモーメントは正であり、極を源へ移し、体が第一座標軸に沿うように軸を回すと、モーメントは消える。二つの読みで $\rho$ の角は一致する。二つのチャート距離が等しく、角が $(0,\pi)$ にあれば、モーメントは $(1-\cos\alpha)\,\chi\,(\chi-\sin\chi\cos\chi)\sin\alpha$ に等しく、正である。そのような源へ平坦な弦に沿って向かう体は、源のまわりの角運動量を得る。その球の外では、内部の配置がどうであれ、速度はバリオン質量全体の点質量の法則になる。重力は半径 \(R\) の銀河パッチの上で閉じる。電磁チャンネルは同じ格子をより大きなコンパクト化半径 \(\Lambda\) で読む。\(\Lambda\) は \(a_0\) と同様、入力である。パッチの内側で重力加速度が逆二乗則を上回る割合は \(x\coth x-1\) であり、\(r=R\) では \(0.31\) と \(0.32\) の間にある。\(r\le R\) の至るところで電磁側の割合は \((R/\Lambda)^2/3\) 未満であり、\(r=R\) では重力側の割合の、電磁側の割合に対する比は \(\frac{93}{100}(\Lambda/R)^2\) を超える。この隔たりは、銀河中心に対して半径方向の、重力チャンネルの向心残差である。

三次元球面には平行線が引ける。純虚の単位四元数は各点で消えない三枠を定め、各方向のノルムは \(1\) で、互いに直交する。左移動 \(t\mapsto q(\cos t+I\sin t)\) に沿う変位ローターは、基点にも始点にも依らず \(\cos s+I\sin s\) である。右移動 \(t\mapsto(\cos t+I\sin t)q\) では同じ変位が左側に立つ。平行線はしたがって二種類あり、いずれも任意の点を通る。共役は左の族を共役点の右の族へ送り、パラメーターの向きを反転する。二種類は鏡像である。恒等元では二つの枠は一致する。測地角 \(\varphi\) の点 \(\cos\varphi+I\sin\varphi\) では、右枠の \(J\) 方向が左枠の \(\cos 2\varphi\, J-\sin 2\varphi\, K\) に重なり、\(0\le\varphi\le\pi/2\) では枠のなす角は \(2\varphi\) である。測地角の向きを返せば余弦は保たれ、\(K\) 成分の符号だけが反転する。角 \(\pi/4\) では二つの \(J\) 方向は一致しない。この倍角が、ホップ・ファイブレーションにツイスター対応が読むねじれである。半径 \(\rho\) の読みでは、物理距離 \(s\) での食い違いは \(2s/\rho\) である。重力チャンネルはこれを半径 \(R\) で読み、電磁チャンネルは半径 \(\Lambda\) で読む。同じ距離での比は \(\Lambda/R\) であり、差は \(2s(\Lambda-R)/(R\Lambda)\) である。パッチの縁 \(s=R\) では重力側の食い違いは \(2\) である。電磁側の半径をハッブル長 \(c/H_0\) と置く読みでは、\(H_0\) を、プランク衛星と距離梯子の中心値をともに含む \(67\) から \(74\,\mathrm{km\,s^{-1}\,Mpc^{-1}}\) の外部区間とし、天の川のコンパクト化半径を \(8\) から \(9\,\mathrm{kpc}\) とすると、縁での電磁側の食い違いは \(3.57\times 10^{-6}\) と \(4.45\times 10^{-6}\) の間にあり、重力側がそれを上回る因子は \(450138\) と \(559315\) の間にある。パッチ半径のハッブル長に対する比はその逆数であり、\(1.787\times 10^{-6}\) と \(2.222\times 10^{-6}\) の間にある。したがって半径 \(\Lambda\) の三次元球面上では、方位等距離図法のヤコビアンはパッチの至るところで \(1\) との差が \(1.646\times 10^{-12}\) 未満であり、電磁側の割合も同じ値を下回る。縁では重力側の割合が電磁側の割合を \(1.8836\times 10^{11}\) より大きく上回る。\(H_0\) は \(a_0\) と同じく入力であり、この同定は双対ローター代数からは出ない。力の割合が \((\Lambda/R)^2\) で隔たる傍らで、二つの平行化そのものは \(\Lambda/R\) で隔たる。三次元球面の回転 \(p\mapsto apb\) は、それぞれの族をそれ自身の上へ運ぶ。右因子は左の線の軸を共役し、左因子は右の線の軸を共役する。族を入れ替え、向きとパラメーターをともに反転するのは、鏡像である。点 \(J\) を通る四分の一回転は、左の族では \(-K\) に、右の族では \(K\) に終わる。\(0<\varphi<\pi/2\) では、枠の \(K\) 成分は前者のねじれで \(-\sin 2\varphi\)、測地角を返したねじれで \(+\sin 2\varphi\) である。どちらも零ではなく、二つは反対であり、回転はそれぞれの族をその族の上に残す。三次元球面が担う族は一つである。そのどちらであるかは、この宇宙の向きである。もう一方は、鏡像の宇宙が担う。

チャート原点にある源は、モーメント \eqref{eq:sphere-chord} を加えない。原点から外れた源は、原点と二つの位置が張る平面の中で、そのモーメントを加える。その残差の上では、円軌道の比角運動量は半径とともに狭義に増し、原点に向かって消えるので、各値はちょうど一つの半径に属する。中心の近くでは、小さな角運動量はなお小さな軌道に属する。コンパクト化半径の三倍を越えると速度は揃い、比角運動量は半径とともに増える。モーメント \eqref{eq:sphere-chord} は、二つの動径の向きのついた角について奇である。源と、半直線を挟んで折り返した源との和は零であり、チャート距離が一致して \((0,\pi)\) にあれば、二つのモーメントはいずれも零ではなく、反対である。この球面のあらゆるパッチは、同じ一つのねじれに対して読まれる。円盤の向きはそのねじれに対する向きであり、基準はどの銀河でも一つである。鏡像の宇宙は、モーメントとともにその基準を反転する。トルクが角運動量を外へ運ぶにつれて、集合体は拡大し、回転は \(v_\infty\) の近くに保たれる。その回転が渦であり、球は円盤へと広がる。

あらゆる半径で
\[
\max\!\left(v_{\text{bar}}^2,\sqrt{a_0\,r\,v_{\text{bar}}^2}\right)<v^2<v_{\text{bar}}^2+\sqrt{a_0\,r\,v_{\text{bar}}^2}
\]
が成り立つ。したがって曲線は、正則な中心のように $r\,v_{\text{bar}}^2$ が消えるところで消える。ケプラー的な尾 $r\,v_{\text{bar}}^2\to GM_{\text{b}}$ をもつあらゆるバリオン曲線は $v_\infty^2=\sqrt{GM_{\text{b}}a_0}$ で平坦になり、バリオンの外では、$R=\sqrt{GM_{\text{b}}/a_0}$ として上の点質量の振る舞いが成り立つ。曲線は立ち上がり、折れ曲がり、平坦なままである。

質量 $6\times10^{10}\,M_\odot$、スケール長 $2.6\,\mathrm{kpc}$ の天の川型の指数円盤に対し、式~\eqref{eq:dual-rotation} は $4\,\mathrm{kpc}$ で約 $202\,\mathrm{km/s}$、$8$ で $211$、$15$ で $190$、$25$ で $179$、$60\,\mathrm{kpc}$ で $176$ を与える。円盤だけのニュートン曲線は、$8$ から $25\,\mathrm{kpc}$ の間に $188$ から $105\,\mathrm{km/s}$ へ落ちる。質量 $3\times10^{9}\,M_\odot$、スケール長 $3\,\mathrm{kpc}$ の低表面輝度円盤は、内側で $g_{\text{N}}$ が $a_0/10$ 程度であり、$2\,\mathrm{kpc}$ で約 $50\,\mathrm{km/s}$、$6$ で $78$、$15\,\mathrm{kpc}$ で $87$ まで上昇し、$83.1$ から $83.2\,\mathrm{km/s}$ の間に落ち着く。質量 $4\times10^{8}\,M_\odot$ で、バリオンが半径 $0.8\,\mathrm{kpc}$ の球の内側にある矮小銀河は、コンパクト化半径が $0.6816$ から $0.6817\,\mathrm{kpc}$ の間にあり、したがって $0.25$ パーセント精度の平坦域は $2.04$ から $2.05\,\mathrm{kpc}$ の間で始まる。漸近速度は $50.2$ から $50.3\,\mathrm{km/s}$ の間にある。球の表面では速度は $55.2$ から $55.3\,\mathrm{km/s}$ の間にあり、$2\,\mathrm{kpc}$ までには $50.38$ から $50.39\,\mathrm{km/s}$ の間まで落ち、$4\,\mathrm{kpc}$ では再び $50.2$ から $50.3\,\mathrm{km/s}$ の間に戻る。その表面から $4\,\mathrm{kpc}$ まで、半径は五倍になり、円軌道の比角運動量は $4.53$ から $4.56$ 倍になる。半径方向の格子であれば、対蹠点は $2.141$ から $2.142\,\mathrm{kpc}$ の間にあり、その平坦な区間の内側にある。

\section{天空が決めうること}
\label{sec:sky}

隣接するパッチは重なりうる。それぞれは自らの中心の近くで忠実である。銀河および銀河団がそのようなパッチであるなら、観測加速度がバリオン的なニュートン値を上回る分は双対像格子であり、見えない質量ではない。直接探索が粒子として見つけるべきものは何もない。三つの特徴が、この読みを暗黒ハローとも、当てはめられた補間とも区別する。バリオンの外で、曲線は上から $v_\infty=(GM_{\text{b}}a_0)^{1/4}$ に近づき、$3R$ を越えると $0.25$ パーセントの精度で平坦である。ケプラー的な尾のように減少することも、再び上昇することもない。バリオン的タリー--フィッシャー関係 $v_\infty^4=GM_{\text{b}}a_0$ は厳密であり、表面輝度にも大きさにも依らない。そして関係 \eqref{eq:dual-rar} は確定した遷移をもち、その遷移は経験的なものより低い。$g_{\text{N}}=a_0$ で増強は $1.31$--$1.32$ であるが、マクゴー、レッリ、ションバートの当てはめはそこで約 $1.58$ を与える。上の銀河では、双対曲線はその当てはめを最大で一割下回り、両者は共通の $v_\infty$ へ向けてゆっくりとしか近づかない。恒星が支配する円盤では、質量光度比をおよそ五分の一高くすればその差は吸収される。バリオン曲線が 21-cm データで固定されるガス支配の矮小銀河では、何もそれを吸収しない。それらが格子と当てはめられた関係との間を決める。強場域における引力と斥力の交替は、等しいラピディティの干渉が秩序づける別のレジームである。このチャートの上の $g$ の符号ではない。円盤の向きを一つのねじれに対して読むことは、同じ天空に、この宇宙が二つの族のどちらを担うかを問う。

\begin{thebibliography}{99}
\raggedright

\bibitem{dst-main}
Anonymous,
「二重時空理論：通常時空と双対時空の間のねじれとしての重力」 (2025).
\url{https://github.com/hypernumbernet/double-spacetime-doc/blob/main/double-spacetime-theory.pdf}.

\end{thebibliography}

\end{document}
